% !TEX spellcheck = en_US

% ~~~~~~~~~~~~~~~~~~~~~~~~~~~~~~~~~~~~~~~ V
% (NC) 2026 Haluk Bingol
% github.com/halukbingol/LaTeX-Templates
%
% Licensees may copy, distribute, display, and perform the work and make derivative works 
% and remixes based on it only for non-commercial purposes. 
% ~~~~~~~~~~~~~~~~~~~~~~~~~~~~~~~~~~~~~~~ A


\documentclass[aspectratio=169,10pt,compress]{beamer}


% ~~~~~~~~~~~~~~~~~~~~~~~~~~~~~~~~~~~~~~~~~~~~~~~~~~~~~~~~~~~~~ V hblib
% Note that hblibCommon is loaded by hblibPresentation
\usepackage{../../cis000_latex/lib/hblibPresentation}
\usepackage{../../cis000_latex/lib/hblibCodeoutput}
\usepackage{../../cis000_latex/lib/hblibDatastructures}
% ~~~~~~~~~~~~~~~~~~~~~~~~~~~~~~~~~~~~~~~~~~~~~~~~~~~~~~~~~~~~~ A hblib



% ~~~~~~~~~~~~~~~~~~~~~~~~~~~~~~~~~~~~~~~ V specific
\usetikzlibrary{fit}
\tikzset{
  var/.style={draw, thick, fill=blue!8, minimum width=1.1cm, minimum height=0.55cm, font=\ttfamily\small},
  obj/.style={draw, thick, rounded corners, fill=orange!10, inner sep=4pt},
  objfit/.style={draw, thick, rounded corners, orange!70!black, inner sep=0pt},
  ref/.style={-{Stealth[length=2.5mm]}, thick, blue!60!black},
  lbl/.style={font=\scriptsize\sffamily, text=gray!70!black}
}
% ~~~~~~~~~~~~~~~~~~~~~~~~~~~~~~~~~~~~~~~ A specific



% ~~~~~~~~~~~~~~~~~~~~~~~~~~~~~~~~~~~~~~~ V title
\newcommand{\cisno}{111}	% 111		
\newcommand{\weekno}{week00}	% week08	
\newcommand{\topic}{ALatexTemplateForPresentation}	% Command Line Interface
\newcommand{\topicshort}{PresentationTemplate}		% CLI
\newcommand{\myauthor}{Haluk Bingol}		% Haluk Bingol
\newcommand{\myauthorshort}{Bingol}		% Bingol
% ~~~~~~~~~~~~~~~~~~~~~~~~~~~~~~~~~~~~~~~ A title

% ~~~~~~~~~~~~~~~~~~~~~~~~~~~~~~~~~~~~~~~ V derived
% do not change
\newcommand{\cisxxx}{\textsc{cis\cisno}}	% do not change
% ~~~~~~~~~~~~~~~~~~~~~~~~~~~~~~~~~~~~~~~ A derived

% ~~~~~~~~~~~~~~~~~~~~~~~~~~~~~~~~~~~~~~~ V titlePage
% Common title format for lab description
% Generates title page
% !TEX spellcheck = en_US

%: ~~~~~~~~~~~~~~~~~~~~~~~~~~~~~~~~~~~~~~ V title
% Common title format for cis presentation

\title[\cisxxx-\weekno-\topicshort]{
	\topic
}

\subtitle{
	\cisxxx-\weekno
}

\author[\myauthorshort]{
	\myauthor
}

\institute[FCIS]{
    Faculty of Computer and Information Sciences\\
    Yeditepe University
}
\date{
	{\tiny \hbTimeStamp}
}
% ~~~~~~~~~~~~~~~~~~~~~~~~~~~~~~~~~~~~~~ A title




\begin{document}




% ~~~~~~~~~~~~~~~~~~~~~~~~~~~~~~~~~~~~~~ V frm
\begin{frame}[t,fragile]
  \titlepage
\end{frame}
% ~~~~~~~~~~~~~~~~~~~~~~~~~~~~~~~~~~~~~~ A frm


% comes ``Outline''
% ~~~~~~~~~~~~~~~~~~~~~~~~~~~~~~~~~~~~~~~ A titlePage

% ~~~~~~~~~~~~~~~~~~~~~~~~~~~~~~~~~~~~~~~ V
% Document starts here
% Outline slide follows this
% ~~~~~~~~~~~~~~~~~~~~~~~~~~~~~~~~~~~~~~~ A




% ~~~~~~~~~~~~~~~~~~~~~~~~~~~~~~~~~~~~~~ V frm
\begin{frame}[t,fragile] \frametitle{Outline}
%  \tableofcontents[hideallsubsections]
  \begin{columns}[T]
    \begin{column}{0.48\textwidth}
      \tableofcontents[sections={1-7}, hideallsubsections]
    \end{column}
    \begin{column}{0.48\textwidth}
      \tableofcontents[sections={8-14}, hideallsubsections]
    \end{column}
  \end{columns}
\end{frame}
% ~~~~~~~~~~~~~~~~~~~~~~~~~~~~~~~~~~~~~~ A frm




% ~~~~~~~~~~~~~~~~~~~~~~~~~~~~~~~~~~~~~~ V
\hSection[Separator]
	{Separators}
% ~~~~~~~~~~~~~~~~~~~~~~~~~~~~~~~~~~~~~~ A




% ~~~~~~~~~~~~~~~~~~~~~~~~~~~~~~~~~~~~~~ V
\hSection[SeparatorSkip]
	{This is ment to be for study at home}
	{\hSectionSkip}
% ~~~~~~~~~~~~~~~~~~~~~~~~~~~~~~~~~~~~~~ A

% ~~~~~~~~~~~~~~~~~~~~~~~~~~~~~~~~~~~~~~ V
\subsection[CodeInText]{Code in text}
% ~~~~~~~~~~~~~~~~~~~~~~~~~~~~~~~~~~~~~~ A




% ~~~~~~~~~~~~~~~~~~~~~~~~~~~~~~~~~~~~~~ V frm
\begin{frame}[t,fragile] \frametitle{Code in text}
  \begin{columns}[T]
    \begin{column}{0.5\textwidth}
    
      \begin{center}
      \begin{tikzpicture}
        \node[var] (a) at (0,0) {a};
        \node[lbl, above=1pt of a] {variable};

        \node[font=\ttfamily\small, anchor=west] at (1.0,0) {null \quad (no object)};
        \node[var] (b) at (0,-1.3) {b};
\node[objfit, minimum width=0.5cm, minimum height=0.5cm] (ob) at (1.65,-1.3) {};

        \node[lbl, right=2pt of ob] {length 0};
        \draw[ref] (b) -- (ob);
        \node[var] (c) at (0,-2.6) {c};
\begin{scope}[shift={(1.4,-2.800)}]\hString[0.4]{\textvisiblespace,\textvisiblespace,\textvisiblespace}\end{scope}
\node[objfit, fit={(1.280,-2.920) (2.720,-2.100)}] (oc) {};

        \node[lbl, right=2pt of oc] {length 3};
        \draw[ref] (c) -- (oc);
      \end{tikzpicture}
      \end{center}


% +++++++++++++++++++++++++++++++++++++++ V lst
%: -Lst.
\begin{lstlisting}[style=codesnippet, language=Java]
String s = "Hello";
int    n = s.length();    // 5
char   c = s.charAt(1);   // 'e'
\end{lstlisting}
% +++++++++++++++++++++++++++++++++++++++ A lst



    \end{column}
    \begin{column}{0.46\textwidth}
      \small
      \begin{tabular}{@{}lccc@{}}
        \toprule
                               & \texttt{a}     & \texttt{b}    & \texttt{c} \\
        \midrule
        \hCode{x == null}     & true           & false         & false \\
        \hCode{x.isEmpty()}   & \emph{error}   & true          & false \\
        \hCode{x.isBlank()}   & \emph{error}   & true          & true \\
        \bottomrule
      \end{tabular}

      \vspace{0.6em}
      
	\begin{alertblock}{NullPointerException}
		\small Calling a method on \hCode{null} fails: there is no object to ask.
		Check \hCode{x != null} first, or write \hCode{"text".equals(x)}.
	\end{alertblock}
	
	\begin{block}{The string pool}
		\small Equal literals can share one object, saving memory (next section).
	\end{block}

    \end{column}
  \end{columns}
\end{frame}
% ~~~~~~~~~~~~~~~~~~~~~~~~~~~~~~~~~~~~~~ A frm


% ~~~~~~~~~~~~~~~~~~~~~~~~~~~~~~~~~~~~~~ V
\hSection[CodeOutput]{
	Code and Output
}
% ~~~~~~~~~~~~~~~~~~~~~~~~~~~~~~~~~~~~~~ A


% ~~~~~~~~~~~~~~~~~~~~~~~~~~~~~~~~~~~~~~ V
\subsection[Right]{Output at the right}
% ~~~~~~~~~~~~~~~~~~~~~~~~~~~~~~~~~~~~~~ A




% ~~~~~~~~~~~~~~~~~~~~~~~~~~~~~~~~~~~~~~ V frm
\begin{frame}[t,fragile] \frametitle{Output at the right}
	\lstcoderightstep[0.75]
		[basicstyle=\ttfamily\scriptsize]
		[basicstyle=\ttfamily\scriptsize]
			{codeoutput/DemoCode.java}{Java}
			{codeoutput/DemoCode.txt}

	\vspace{0.3em}
	Some text
\end{frame}
% ~~~~~~~~~~~~~~~~~~~~~~~~~~~~~~~~~~~~~~ A frm



% ~~~~~~~~~~~~~~~~~~~~~~~~~~~~~~~~~~~~~~ V
\subsection[on]{Output on code}
% ~~~~~~~~~~~~~~~~~~~~~~~~~~~~~~~~~~~~~~ A




% ~~~~~~~~~~~~~~~~~~~~~~~~~~~~~~~~~~~~~~ V frm
\begin{frame}[t,fragile] \frametitle{Output on code}

	\lstcodeoverlaystep
		[basicstyle=\ttfamily\scriptsize]
		[basicstyle=\ttfamily\scriptsize][0.28\textwidth]{12.0cm,2.25cm}
			{codeoutput/DemoCode.java}{Java}
			{codeoutput/DemoCode.txt}

	\vspace{0.2em}
	Some text
\end{frame}
% ~~~~~~~~~~~~~~~~~~~~~~~~~~~~~~~~~~~~~~ A frm




% ~~~~~~~~~~~~~~~~~~~~~~~~~~~~~~~~~~~~~~ V
\subsection[bottom]{Output at the bottom}
% ~~~~~~~~~~~~~~~~~~~~~~~~~~~~~~~~~~~~~~ A




% ~~~~~~~~~~~~~~~~~~~~~~~~~~~~~~~~~~~~~~ V frm
\begin{frame}[t,fragile] \frametitle{Output at the bottom}
	\lstcodedownstep
		[basicstyle=\ttfamily\scriptsize]
		[basicstyle=\ttfamily\scriptsize, linerange={1-2,13-13}]
			{codeoutput/DemoCode.java}{Java}
			{codeoutput/DemoCode.txt}
	Some text
\end{frame}
% ~~~~~~~~~~~~~~~~~~~~~~~~~~~~~~~~~~~~~~ A frm



% ~~~~~~~~~~~~~~~~~~~~~~~~~~~~~~~~~~~~~~ V
\section{Algorithm}
% ~~~~~~~~~~~~~~~~~~~~~~~~~~~~~~~~~~~~~~ A



% ~~~~~~~~~~~~~~~~~~~~~~~~~~~~~~~~~~~~~~ V
\subsection[Alg]{Algorithm}
% ~~~~~~~~~~~~~~~~~~~~~~~~~~~~~~~~~~~~~~ A




% ~~~~~~~~~~~~~~~~~~~~~~~~~~~~~~~~~~~~~~ V frm
\begin{frame}[t,fragile] \frametitle{Algorithm}
  \begin{columns}[T]
    \begin{column}{0.5\textwidth}
      \small A \textbf{palindrome} reads the same backwards.\\
      Examples:\\
	\begin{itemize}
		
		\item 
		\texttt{kayak},
		
		\item 
		\texttt{level},
		
		\item 
		\texttt{ey edip adanada pide ye}.
		
	\end{itemize}
      

      \vspace{0.4em}
      Compare the two ends, then move inwards:





	% +++++++++++++++++++++++++++++++++++++++ V
	%: -Fig.      
      \begin{center}
      \begin{tikzpicture}
        \hString[0.55]{k,a,y,a,k}
        \draw[{Stealth}-{Stealth}, thick, orange] (0.27,-0.1) to[bend right=35] (2.47,-0.1);
        \draw[{Stealth}-{Stealth}, thick, blue] (0.82,-0.1) to[bend right=35] (1.92,-0.1);
        \node[font=\scriptsize] at (1.37,-0.95) {\texttt{i} moves right, \texttt{j} moves left};
      \end{tikzpicture}
      \end{center}
	% +++++++++++++++++++++++++++++++++++++++ A
    \end{column}
    \begin{column}{0.46\textwidth}
      \SetAlFnt{\small}


% +++++++++++++++++++++++++++++++++++++++ V alg
%: -Alg. 
\begin{algorithm}[H]
  \caption{IsPalindrome($s$)}
  $i \leftarrow 0$, \quad $j \leftarrow \mathit{length}(s) - 1$\;
  \While{$i < j$}{
    \If{$s[i] \neq s[j]$}{
      \Return{false}\;
    }
    $i \leftarrow i + 1$, \quad $j \leftarrow j - 1$\;
  }
  \Return{true}\;
\end{algorithm}
% +++++++++++++++++++++++++++++++++++++++ A alg

    \end{column}
    \begin{column}{0.46\textwidth}

    \end{column}
  \end{columns}
\end{frame}
% ~~~~~~~~~~~~~~~~~~~~~~~~~~~~~~~~~~~~~~ A frm




% ~~~~~~~~~~~~~~~~~~~~~~~~~~~~~~~~~~~~~~ V
\hSection[References]
	{References}
% ~~~~~~~~~~~~~~~~~~~~~~~~~~~~~~~~~~~~~~ A






% ~~~~~~~~~~~~~~~~~~~~~~~~~~~~~~~~~~~~~~ V frm
\begin{frame}[t,fragile] \frametitle{References}
  \begin{columns}[T]
  	How to cite
    \begin{column}{0.5\textwidth}
	01pi.com~\cite{pi10ObjectsComplex}
    \end{column}
    \begin{column}{0.5\textwidth}
	CLRS~\cite{cormen2022introduction}
    \end{column}
  \end{columns}
\end{frame}
% ~~~~~~~~~~~~~~~~~~~~~~~~~~~~~~~~~~~~~~ A frm


% ~~~~~~~~~~~~~~~~~~~~~~~~~~~~~~~~~~~~~~ V References bib
\begin{frame}[t,allowframebreaks]\frametitle{References}
	\setbeamertemplate{bibliography item}[text] % sets numbers in references
%	\bibliographystyle{abbrv}
%	\bibliographystyle{apalike}
%	\bibliographystyle{ieeetr}
	\bibliographystyle{IEEEtran}
	%
	\tiny\bibliography{../../cis000_latex/lib/hblibDataStructures} % bib file

	
\end{frame}
% ~~~~~~~~~~~~~~~~~~~~~~~~~~~~~~~~~~~~~~ A





% ~~~~~~~~~~~~~~~~~~~~~~~~~~~~~~~~~~~~~~ V
\hSection[Summary]
	{Summary}
% ~~~~~~~~~~~~~~~~~~~~~~~~~~~~~~~~~~~~~~ A




% ~~~~~~~~~~~~~~~~~~~~~~~~~~~~~~~~~~~~~~ V frm
\begin{frame}[t,fragile]
	\begin{center}
		\vspace{4cm}
		\hRemark{\Huge Questions?}
	\end{center}
\end{frame}
% ~~~~~~~~~~~~~~~~~~~~~~~~~~~~~~~~~~~~~~ A frm



\end{document}
